
\documentclass[11pt]{article}
\usepackage{amsmath,amssymb,amsfonts}
\usepackage{geometry}
\geometry{margin=1in}

\title{CEQG-RG-Langevin with Multiplicity-Native Mathematically Inverted Digital Twin\\\\Tracks A, B, C: Comprehensive Mathematical Overview}
\author{Ryan O. Van Gelder \\ \\ Citizen Gardens}
\date{March 2026}

\begin{document}
\maketitle

\section{Introduction}

This report summarizes the current status of the CEQG-RG-Langevin framework augmented with three parallel UV/microscopic tracks:

\begin{itemize}
  \item Track A: Pure Arithmetic-Langevin Group Field Theory (AL-GFT) based on EPRL spin foams.
  \item Track B: String-enhanced module, in which a 2D stringy worldsheet sector acts as an additional environmental bath (Mikovi\'c-type string spin foam and string-field-theory motivated RG priors).
  \item Track C: Multiplicity-native, mathematically inverted digital twin, defined as the inverse problem of the forward CEQG pipeline, constructed directly from the prime-labeled multiplicity cumulant hierarchy.
\end{itemize}

All three tracks feed a unified Einstein--Langevin backbone, share the same five-gate falsifiability structure, and differ only in their microscopic/environmental realization and corresponding renormalization-group (RG) priors.

\section{Unified Einstein--Langevin Backbone}

Let $g_{\mu\nu}$ denote the spacetime metric and $G_{\mu\nu}[g]$ the Einstein tensor. The coarse-grained semi-classical dynamics, including backreaction from quantum fluctuations of matter and geometry, is encoded in the Einstein--Langevin equation
\begin{equation}
  G_{\mu\nu}[g] + \Lambda_{\rm eff}(M) g_{\mu\nu}
  = 8\pi G \big( \langle T_{\mu\nu} \rangle_{\rm ren} + \xi^{(X)}_{\mu\nu} \big),
  \label{eq:einstein-langevin}
\end{equation}
where $\Lambda_{\rm eff}(M)$ is an effective scale-dependent cosmological term, $\langle T_{\mu\nu} \rangle_{\rm ren}$ is the renormalized expectation value of the stress--energy tensor, and $\xi^{(X)}_{\mu\nu}$ is a stochastic tensor encoding fluctuations, with $X \in \{A,B,C\}$ labelling the track.

The statistics of $\xi^{(X)}_{\mu\nu}$ are determined by a noise kernel
\begin{equation}
  \mathcal{N}^{(X)}_{\mu\nu\alpha\beta}(x,x')
  = \frac{1}{2}
    \big\langle \{\hat{t}_{\mu\nu}(x),\hat{t}_{\alpha\beta}(x')\} \big\rangle_{\mathcal{H}_{\rm env}^{(X)}},
  \label{eq:noise-kernel}
\end{equation}
where $\hat{t}_{\mu\nu} = \hat{T}_{\mu\nu} - \langle \hat{T}_{\mu\nu}\rangle$ and $\mathcal{H}_{\rm env}^{(X)}$ is the environmental Hilbert space attached to track $X$.

The environment enters through a Schwinger--Keldysh influence functional $\mathcal{F}^{(X)}[g_{+},g_{-}]$; differentiating $\mathcal{F}^{(X)}$ yields both the dissipation kernel and the noise kernel~\eqref{eq:noise-kernel} via a generalized fluctuation--dissipation relation.

\section{Track A: Pure AL--GFT (Spin-Foam Arithmetic--Langevin)}

\subsection{Microscopic structure}

Track A is based on an EPRL-type spin-foam Group Field Theory (GFT). The microscopic degrees of freedom are group-valued fields
\begin{equation}
  \varphi(g_1,\dots,g_d): G^d \to \mathbb{C},
\end{equation}
with $d=4$ for 4D gravity and $G=\mathrm{SU}(2)$ (or a suitable extension). The GFT action reads schematically
\begin{equation}
  S_{\rm GFT}[\varphi,\bar{\varphi}] = S_{\rm kin} + S_{\rm int},
\end{equation}
where $S_{\rm kin}$ encodes propagation on $G^d$ and $S_{\rm int}$ reproduces EPRL spin-foam amplitudes in perturbation theory. The partition function is
\begin{equation}
  \mathcal{Z}_{\rm GFT} = \int \mathcal{D}\varphi\,\mathcal{D}\bar{\varphi}\, e^{-S_{\rm GFT}[\varphi,\bar{\varphi}]}.
\end{equation}

\subsection{Arithmetic--Langevin environment and zeta comb}

The environmental modes are encoded as a set of harmonic degrees of freedom $\{\phi_{n,k}\}$ with frequencies $\omega_n$ and couplings
\begin{equation}
  g_n = \varepsilon \, e^{-\gamma \omega_n^2} e^{i\phi_n},
\end{equation}
where the phases $\phi_n$ carry an arithmetic imprint (e.g. related to non-trivial zeros of the Riemann zeta function). Integrating out these modes along the Schwinger--Keldysh contour yields an influence functional $\mathcal{F}^{(A)}[g_{+},g_{-}]$ and a corresponding noise kernel $\mathcal{N}^{(A)}_{\mu\nu\alpha\beta}$.

\subsection{RG flow and Wetterich equation}

The effective average action $\Gamma_k$ obeys a Wetterich-type RG equation
\begin{equation}
  \partial_k \Gamma_k[\Phi]
  = \frac{1}{2} \mathrm{Tr}
    \Big\{
      (\Gamma_k^{(2)}[\Phi] + R_k)^{-1} \partial_k R_k
    \Big\},
\end{equation}
with truncations chosen to respect GFT power counting, Ward identities and EPRL consistency. In particular, the scale dependence of effective couplings such as $\lambda_4(k)$, $\lambda_6(k)$ and a log-link parameter $\nu_{\rm eff}(M)$ is determined within this flow.

\section{Track B: String-Enhanced Environmental Sector}

\subsection{Worldsheet sector and stringy bath}

Track B augments the AL--GFT environment with a 2D stringy worldsheet sector. The total environmental Hilbert space factorizes as
\begin{equation}
  \mathcal{H}_{\rm env}^{(B)} = \mathcal{H}_{\rm ws} \otimes \mathcal{F}_{\rm GFT},
\end{equation}
where $\mathcal{H}_{\rm ws}$ is the Fock space of bosonic string oscillators on a group manifold and $\mathcal{F}_{\rm GFT}$ is the Fock space of GFT quanta (as in Track A).

The worldsheet partition function on a hyperbolic surface $\Gamma\backslash\mathbb{H}$ can be expressed in terms of a Selberg zeta function $Z_\Gamma(s)$. The resulting contribution to the noise kernel can be written schematically as a spectral correction
\begin{equation}
  \delta \mathcal{N}^{(\mathrm{ws})}(\omega)
  = \alpha_{\rm ws}\; \mathrm{Re}\left[\frac{Z_\Gamma'(\tfrac{1}{2} + i\omega)}{Z_\Gamma(\tfrac{1}{2} + i\omega)}\right],
\end{equation}
with a coupling $\alpha_{\rm ws}$ controlling the strength of the worldsheet contribution. The full Track-B noise kernel is then
\begin{equation}
  \mathcal{N}^{(B)} = \mathcal{N}^{(A)} + \delta \mathcal{N}^{(\mathrm{ws})}.
\end{equation}

\subsection{String-motivated RG priors and Swampland flag}

The RG flow of effective couplings in Track B receives additional contributions from worldsheet beta functions. For an effective parameter $\nu_{\rm eff}(M)$ one can write
\begin{equation}
  M\partial_M \nu_{\rm eff}^{(B)}
  = \beta_{\rm GFT}(\nu_{\rm eff}) + \beta_{\rm ws}(\nu_{\rm eff}, \alpha_{\rm ws}),
\end{equation}
where $\beta_{\rm GFT}$ is the Track-A contribution and $\beta_{\rm ws}$ encodes stringy corrections. Swampland-inspired constraints (distance conjecture, de Sitter conjecture, etc.) may be implemented as a conditional prior with a binary flag $\delta_{\rm sw} \in \{0,1\}$, distinguishing runs with and without these conjectural constraints active.

\section{Track C: Mathematically Inverted Digital Twin}

Track C is defined as the multiplicity-native inverse problem of the forward CEQG pipeline. Tracks A and B proceed from microscopic spin-foam amplitudes $A_{\rm micro}$ (EPRL or string-enhanced) through coarse-graining (Schwinger--Keldysh influence functional plus GFT/Wetterich flow) to macroscopic noise kernels and cumulants. Track C instead starts from a prescribed target macroscopic multiplicity structure $\{\kappa_n^{\rm MT}\}$ (the hierarchy of connected cumulants fixed by the prime-labeled recursion of Multiplicity Theory) and solves the inverse map to reconstruct a consistent microscopic measure $\mu_{\rm MT}$.

\subsection{Inverse problem formulation}

Let $C_n^{\rm IR}[\mu]$ denote the IR cumulants obtained by running the forward CEQG pipeline (influence functional + RG flow) starting from a microscopic measure $\mu$. Track C defines $\mu_{\rm MT}$ as the solution of the inverse problem
\begin{equation}
  \mu_{\rm MT}
  =
  \arg\inf_{\mu}
  \Bigl[
    D\bigl(C_n^{\rm IR}[\mu] \,\Vert\, \kappa_n^{\rm MT}\bigr)
    + \lambda\,\mathrm{Reg}(\mu)
  \Bigr],
  \label{eq:inverse-problem}
\end{equation}
where $D$ is a statistical divergence (e.g. Kullback--Leibler) between the IR cumulants generated from $\mu$ and the target multiplicity recursion $\kappa_n^{\rm MT}$, $\mathrm{Reg}(\mu)$ is a regularization term enforcing GFT power counting, Ward identities, and EPRL-consistency constraints, and $\lambda$ is a hyperparameter tuned at Gate 2.

From the reconstructed $\mu_{\rm MT}$ one extracts the implied micro-amplitudes $A_{\rm MT}$ and the corresponding influence functional $\mathcal{F}_{\rm MT}[\phi]$, yielding a Track-C noise kernel $\mathcal{N}^{(C)}_{\mu\nu\alpha\beta}$ via Eq.~\eqref{eq:noise-kernel} with $\mathcal{H}_{\rm env}^{(C)}$ generated from $\mu_{\rm MT}$.

\subsection{Forward vs inverse duality}

Equation~\eqref{eq:inverse-problem} makes Track C a mathematically inverted digital twin of the forward Tracks A/B: instead of starting from $A_{\rm micro}$ and predicting $\{C_n^{\rm IR}\}$, Track C starts from $\{\kappa_n^{\rm MT}\}$ and infers a compatible $\mu_{\rm MT}$ (and thus $A_{\rm MT}$). This realizes an explicit forward/inverse duality at the level of the full CEQG-RG-Langevin pipeline.

\section{Unified Noise Kernel and Bayesian Evidence}

Across the three tracks, the noise kernel entering the Einstein--Langevin equation~\eqref{eq:einstein-langevin} has the unified form
\begin{equation}
  \mathcal{N}^{(X)}_{\mu\nu\alpha\beta}(x,x')
  =
  \frac{1}{2}
  \big\langle
    \{\hat{t}_{\mu\nu}(x), \hat{t}_{\alpha\beta}(x')\}
  \big\rangle_{\mathcal{H}_{\rm env}^{(X)}},
  \qquad X \in \{A,B,C\}.
\end{equation}

Given observational data $\mathcal{D}$ (e.g. CMB spectra, non-Gaussianity running, gravitational-wave propagation), each track defines a likelihood $p(\mathcal{D}\mid X,\mathcal{M})$ under the shared model architecture $\mathcal{M}$. The Bayes factors
\begin{equation}
  \mathcal{K}(B\mid A) = \frac{p(\mathcal{D}\mid B,\mathcal{M})}{p(\mathcal{D}\mid A,\mathcal{M})},
  \qquad
  \mathcal{K}(C\mid A) = \frac{p(\mathcal{D}\mid C,\mathcal{M})}{p(\mathcal{D}\mid A,\mathcal{M})},
  \qquad
  \mathcal{K}(C\mid B) = \frac{p(\mathcal{D}\mid C,\mathcal{M})}{p(\mathcal{D}\mid B,\mathcal{M})},
\end{equation}
quantify whether the multiplicity structure is better realized as an emergent property of forward tracks (A/B) or as a fundamental input defining the inverted track (C).

In practice, these Bayes factors can be approximated via differences in effective chi-squared and Wetterich-derived free energies,
\begin{equation}
  \mathcal{K}(Y\mid X)
  \approx
  \exp\Bigl(-\tfrac{1}{2}\Delta\chi^2_{YX} + \Delta\ln\mathcal{Z}_{YX}\Bigr),
\end{equation}
where $\Delta\chi^2_{YX}$ encodes the fit to data and $\Delta\ln\mathcal{Z}_{YX}$ the difference in RG-improved partition functions between tracks $X$ and $Y$.

\section{Outlook}

The three-track CEQG-RG-Langevin architecture described here is designed to be modular and falsifiable. Track A (pure AL--GFT) provides a baseline spin-foam realization, Track B (string-enhanced) introduces a worldsheet-modular correction layer with Swampland-conditional priors, and Track C (mathematically inverted digital twin) implements a multiplicity-native inverse problem that treats the prime-labeled cumulant hierarchy as primary data.

Numerical implementation of the Track-B worldsheet correction and the Track-C inverse solver (e.g. via adjoint methods for the Wetterich flow) will enable a three-way Bayesian race, with forthcoming observations (e.g. LiteBIRD, DESI) determining whether the multiplicity structure manifests most naturally as emergent, string-modular, or fundamental.

\end{document}